% Created 2023-12-27 mié 22:52
% Intended LaTeX compiler: pdflatex
\documentclass[11pt]{article}
\usepackage[utf8]{inputenc}
\usepackage[T1]{fontenc}
\usepackage{graphicx}
\usepackage{longtable}
\usepackage{wrapfig}
\usepackage{rotating}
\usepackage[normalem]{ulem}
\usepackage{amsmath}
\usepackage{amssymb}
\usepackage{capt-of}
\usepackage{hyperref}
\usepackage{minted}

\usepackage{parskip}
\usepackage[utf8]{inputenc} %% For unicode chars
\usepackage{placeins}
\usepackage[margin=2.50cm]{geometry}
\usepackage[T1]{fontenc}
\usepackage{mathpazo}
\linespread{1.05}
\usepackage[scaled]{helvet}
\usepackage{courier}
\hypersetup{colorlinks=true,linkcolor=blue}
\RequirePackage{fancyvrb}
\DefineVerbatimEnvironment{verbatim}{Verbatim}{fontsize=\small,formatcom = {\color[rgb]{0.5,0,0}}}
\usepackage{mdframed}
\BeforeBeginEnvironment{minted}{\begin{mdframed}}
\AfterEndEnvironment{minted}{\end{mdframed}}
\setcounter{secnumdepth}{3}
\date{}
\title{}
\hypersetup{
 pdfauthor={},
 pdftitle={},
 pdfkeywords={},
 pdfsubject={},
 pdfcreator={Emacs 28.1 (Org mode 9.6.12)}, 
 pdflang={Spanish}}
\begin{document}

\setcounter{tocdepth}{3}
\tableofcontents

\setlength\parindent{10pt}


\section{PHP and HTML}
\label{sec:orgd224f90}

Integrating PHP with HTML and managing data formats are crucial skills
in web development. Here's a suggested order for tackling these topics,
along with explanations and examples for each:

\subsection{Embedding PHP in HTML}
\label{sec:orgc0505e7}
Start by learning how to embed PHP code within an HTML document. This is
foundational for creating dynamic web pages.

\subsubsection{Example:}
\label{sec:orgd2054ac}
\begin{minted}[]{html}
<!DOCTYPE html>
<html>
<head>
    <title>PHP in HTML Example</title>
</head>
<body>
    <h1>Welcome to My Website</h1>
    <?php
        echo "<p>Today's date is " . date("Y-m-d") . ".</p>";
    ?>
</body>
</html>
\end{minted}

This example demonstrates how to insert PHP within an HTML document to
display the current date.

\subsection{Embedding HTML in PHP}
\label{sec:org71fefc5}
Next, understand how to output HTML content from PHP scripts. This is
useful when you need to generate complex HTML structures based on
certain conditions or data processing in PHP.

\subsubsection{Example:}
\label{sec:org72e313a}
\begin{minted}[]{php}
<?php
$loggedIn = true;
echo "<div>";
if ($loggedIn) {
    echo "<h2>Welcome back, user!</h2>";
} else {
    echo "<h2>Welcome, guest!</h2>";
}
echo "</div>";
?>
\end{minted}

In this example, PHP decides which HTML content to output based on a
condition.

\subsection{Using PHP and HTML with Forms}
\label{sec:orgeec87f7}
Learn to create HTML forms and process form data with PHP. This is
essential for user interaction on websites.

\subsubsection{Example:}
\label{sec:org1156daa}
\begin{minted}[]{html}
<!-- HTML Form -->
<form action="submit.php" method="post">
    Name: <input type="text" name="name">
    <input type="submit">
</form>
\end{minted}

\begin{minted}[]{php}
<?php
// submit.php
if ($_SERVER["REQUEST_METHOD"] == "POST") {
    echo "<p>Hello, " . htmlspecialchars($_POST['name']) . "!</p>";
}
?>
\end{minted}

This demonstrates a simple form submission and data handling in PHP.

\subsection{PHP for Control Structures in HTML}
\label{sec:orgcb0113f}
Utilize PHP's control structures (loops, conditionals) to dynamically
generate HTML content.

\subsubsection{Example:}
\label{sec:org22696a2}
\begin{minted}[]{php}
<ul>
    <?php
    for ($i = 1; $i <= 5; $i++) {
        echo "<li>Item $i</li>";
    }
    ?>
</ul>
\end{minted}

This example uses a PHP loop to create a list in HTML.

\subsection{Data Formatting and Manipulation in PHP}
\label{sec:orge5820d0}
Once comfortable with embedding PHP in HTML and vice versa, focus on
data formatting in PHP. This includes string formatting, date/time
formatting, number formatting, etc.

\subsubsection{Example:}
\label{sec:org50a2b01}
\begin{minted}[]{php}
<?php
$price = 49.99;
$formattedPrice = number_format($price, 2);
echo "<p>The price is $$formattedPrice</p>";
?>
\end{minted}

This illustrates number formatting in PHP.

\subsection{Advanced Data Formatting}
\label{sec:org2811f27}
Delve into more complex data formatting like JSON encoding/decoding,
which is particularly useful in web APIs.

\subsubsection{Example:}
\label{sec:org32337ec}
\begin{minted}[]{php}
<?php
$data = ['name' => 'John', 'age' => 30];
$jsonData = json_encode($data);
echo $jsonData;
?>
\end{minted}

This example converts a PHP array into a JSON string.

\subsection{Conclusion and Order of Learning}
\label{sec:org1bd0311}
Starting with embedding PHP in HTML helps set the stage for dynamic web
content creation. As you progress, learning to process and handle user
input via forms becomes essential. Then, focusing on PHP's control
structures within HTML will deepen your understanding of creating
responsive web pages. Finally, mastering data formatting ensures that
the data you work with is presented and utilized effectively.

By following this progression, you'll develop a strong foundation in PHP
for web development, paving the way for more advanced topics such as
database integration, session management, and security.
\end{document}
